% Preface: the README's "Om kursen", "Efter kursen" and "Att bygga", info/README.md's
% prerequisites and software list, and the course's three verification layers, as the front
% matter of a book.

\chapter{Preface}
\markboth{Preface}{Preface}

This book takes you from a gate network on paper to a CAN node in an FPGA that a microcontroller
can talk to over SPI.

The first half builds up VHDL and synchronous design from the ground: combinational networks,
sequential networks, metastability and synchronisation, counters, shift registers, timers and
state machines, each one first reasoned about as a \emph{circuit} and only then expressed in the
language. The second half is a group project where that knowledge is put to something larger: a
simplified CAN controller, designed module by module, and the two layers that make it reachable
from the outside \textemdash{} a register bank and an SPI bridge.

\textbf{The subject is VHDL and digital design in an FPGA, not the fundamentals of digital
electronics.} Gates, truth tables, Boolean algebra and Karnaugh maps are revised rather than
taught, and they are revised where the language calls for them. What the book teaches from zero is
the language and the toolchain.

\section*{What you are assumed to bring}
\begin{itemize}
  \item Logic gates and truth tables, and reading a sum of products out of a table.
  \item Boolean algebra: the basic identities and De Morgan's laws.
  \item Karnaugh maps, for functions of three or four variables.
  \item Registers, D flip-flops, and what a clock edge does.
  \item State machines as a concept: states, transitions, and the difference between Moore and
    Mealy.
\end{itemize}

If any of it needs brushing up, all of it is in \debook{}
(\url{https://github.com/qrtech-academy/digital-electronics}), a freely available book that goes
from a single gate to a state machine: gates and Boolean algebra in chapters~1 and~2, Karnaugh maps
in chapter~3, flip-flops and registers in chapter~5 and state machines in chapter~6. The chapters of
this book point there wherever they revise the basics, and \secref{app:reading:basics} has the whole
map.

No programming language beyond general software experience is assumed. Familiarity with C or
something like it helps, since several explanations set VHDL's semantics against it, and the one
section that actually contains C++ (\secref{c10:sec:framing}) is a short introduction to framing
rather than an exercise in the language.

The group project adds \textbf{Git and GitHub} at a basic level: clone, branch, commit, push and
pull request.

\section*{What you should be able to do afterwards}
\begin{itemize}
  \item Realise a Boolean function as a gate network, and simulate it by hand.
  \item Read and write combinational and sequential VHDL: entities, architectures, signals,
    variables and processes.
  \item Explain metastability and apply the two-flip-flop synchroniser to every asynchronous
    input.
  \item Design counters, shift registers, timers and state machines, and compose a design out of
    modules you have already built rather than rewriting their logic.
  \item Design the hardware of a communication protocol: bit timing, checksums, bit stuffing,
    frame sequencing and arbitration.
  \item Design the interface between hardware and software: a register bank and an SPI transport,
    against a contract someone else writes to.
  \item Verify a design by running a self-checking testbench under GHDL, and trace a failure back
    to the code that caused it.
  \item Work in a group with feature branches, pull requests and code review.
\end{itemize}

\section*{How the book is laid out}
The book is typeset from a course of twenty sessions, and \textbf{every chapter is one session}.
The chapter's sections are the session's appendices, in the same order and under the same
subheadings, so that a chapter here and a lecture directory in the course repository can be read
side by side. The two parts match the course's two halves:

\begin{booktable}{@{}p{4.6em}p{3.6em}L@{}}
  \thead{Chapter} & \thead{Session} & \thead{Subject} \\
  \midrule
  \multicolumn{3}{@{}l}{\textit{Part \ref{part:vhdl}: digital design in VHDL}} \\
  1 & \file{L01} & Combinational logic and first VHDL \\
  2 & \file{L02} & Larger networks, multiplexers and submodules \\
  3 & \file{L03} & Sequential networks \\
  4 & \file{L04} & Metastability and synchronisation \\
  5 & \file{L05} & Variables and the hardware underneath \\
  6 & \file{L06} & Counters and shift registers \\
  7 & \file{L07} & Timers \\
  8 & \file{L08} & State machines \\
  9 & \file{L09} & Practical exam 1: sequential networks \\
  \addlinespace
  \multicolumn{3}{@{}l}{\textit{Part \ref{part:can}: the group project, a CAN node}} \\
  10 & \file{L10} & Project start: the CAN bus, the frame and \code{can_def} \\
  11 & \file{L11} & Architecture, the top level and the register map \\
  12 & \file{L12} & Synchronisation and the bit timer \\
  13 & \file{L13} & The CRC-15 engine \\
  14 & \file{L14} & Transmit shift register and bit stuffing \\
  15 & \file{L15} & Receive shift register and destuffing \\
  16 & \file{L16} & The CAN controller I: the state machine and the transmit path \\
  17 & \file{L17} & The CAN controller II: reception and arbitration \\
  18 & \file{L18} & The register bank and SPI from the waveform \\
  19 & \file{L19} & The SPI bridge, \code{can_spi_node} and bring-up \\
  20 & \file{L20} & Practical exam 2: CAN modules \\
\end{booktable}

Every chapter opens with what is in it, works its way through the material, closes with a review
of what you should be able to explain afterwards, and ends with its exercises. The appendices hold
the tool references, the project specification, the shared contract with the driver course, and a
closing note on where to read next.

Chapters~\ref{c9:ch} and \ref{c20:ch} are the two examinations. They contain no new theory; they
say what the exam measures, what shape the tasks take and how to prepare, and they are worth
reading in advance rather than the day before.

\textbf{One more book belongs to part~\ref{part:can}.} This book describes how to \emph{build} a CAN
node; the protocol itself is described in \canbook{}, a short book written for both halves of a CAN
system, in English and in Swedish with the same chapter numbering. It is the protocol reference
throughout the second half: its chapters 1--6 are the requirement specification for what is built
here, bit by bit, and its chapter~7 describes exactly the simplified controller part~\ref{part:can}
arrives at. Every chapter that adds a protocol rule says where in that book the rule is stated, and
\secref{app:reading:protocol} has the full chapter map between the two. Its chapters 1 and 2 are
worth reading before \chapref{c10:ch}; the rest works best as a reference.

\section*{Three layers of verification}
The course checks its work at three levels, and keeping them apart is a point in itself:

\begin{booktable}{@{}p{9.2em}L@{}}
  \thead{Layer} & \thead{What it checks} \\
  \midrule
  \textbf{By hand} & In CircuitVerse, before any VHDL exists. You derive what the circuit should
    do, build it, and simulate it against your own prediction. Nothing checks that step
    automatically, deliberately: predicting a behaviour and then confirming it is the skill the
    rest of the book rests on. \\
  \textbf{In the exercises} & You write the VHDL and run a testbench handed out with it,
    self-checking, under GHDL. No FPGA board is needed, and nearly everything can be verified on
    an ordinary laptop. \\
  \textbf{On the board} & The design is synthesised in Quartus Prime Lite and run on a Terasic
    DE0-CV. During the first half that step is demonstrated; in the group project's last session
    you do it yourself. \\
\end{booktable}

\Secref{c2:sec:testbenches} is the full introduction to testbenches, and appendix~\ref{app:sim} is
the reference you will come back to. Appendix~\ref{app:quartus} is the Quartus flow.

\section*{What the exercises look like}
The exercises come after every chapter, numbered within the chapter, and each one is tagged with
its kind:

\begin{booktable}{@{}p{7.4em}L@{}}
  \thead{Kind} & \thead{What it asks of you} \\
  \midrule
  \kindtag[fillaccent5]{Circuit} & Draw, build and simulate a circuit by hand in CircuitVerse, and
    check it against your own truth table or your own timing diagram. \\
  \kindtag[fillaccent]{VHDL} & Write a module, with the entity name and the \textbf{port order}
    the exercise states, and run the testbench handed out with it. \\
  \kindtag[fillaccent3]{Simulation} & Change something and find out what GHDL reports. Several of
    them ask you to \emph{predict} the outcome first and then run. \\
  \kindtag[fillaccent2]{Reasoning} & Answer in prose or in numbers: what the code does, why one
    construction is better than another, or what a number works out to. \\
  \kindtag[fillaccent4]{C++} & The course's only C++ exercise, in \chapref{c10:ch}. \\
\end{booktable}

\textbf{There are no printed solutions, and that is deliberate.} Every exercise that asks you to
write a VHDL module comes with a self-checking testbench, and that testbench is the answer key: it
drives your module through every case and says exactly which check failed. The pen-and-paper
exercises you check against the truth table or the timing diagram you derived yourself, which is
the whole point of them. A list of answers would remove precisely the step the exercise exists
for.

\begin{aside}
\textbf{The port order is the contract.} Every testbench handed out binds to the module it drives
\textbf{positionally}, and so does every instantiation in the course. Nothing binds a port by
name. The contract is therefore the \emph{order} and the \emph{types} of the ports, exactly as
every exercise's interface table lists them, top to bottom. The \emph{names} of the ports are
yours.

That bargain is worth understanding rather than merely following. Positional binding is shorter to
write and to read, and it makes the interface table the single source of truth. What it gives up
is the compiler's help: swap two ports of the same type and nothing complains during analysis, the
design simply misbehaves in simulation instead. The book keeps its port names all the same, so that
the text and your code speak about the same signals.
\end{aside}

\section*{Setting up the tools}
Everything in this book runs in a Linux terminal: on Linux, or on Windows through WSL. Three
things cover everything that is not a board:

\begin{itemize}
  \item \textbf{GHDL}, the open-source VHDL analyser and simulator. It is the tool the course runs
    on: every module is verified against a testbench in GHDL, and nothing you are asked to produce
    needs anything else.
  \item \textbf{CircuitVerse}, a free logic simulator that runs in the browser
    (\url{https://circuitverse.org/simulator}). Nothing to install.
  \item \textbf{Git}, for the group project, plus \code{g++} for \chapref{c10:ch}'s framing
    exercise.
\end{itemize}

\noindent On WSL or Ubuntu:

\begin{shell}
sudo apt -y update
sudo apt -y install git make ghdl g++
ghdl --version           # 3.x or later
\end{shell}

\noindent The course repository has the test framework as a submodule, so clone it with the
submodules:

\begin{shell}
git clone --recurse-submodules <repo-url>
git submodule update --init       # if it is already cloned without them
\end{shell}

\noindent Two commands at the root of the repository cover most of what the book asks for:

\begin{shell}
make build               # Build and simulate everything there is: examples, exercises, project.
make lint                # Markdown links, duplicate check, no trailing whitespace.
\end{shell}

\noindent\code{make build-project} reports every project testbench as skipped in a freshly checked
out repository. That is the expected state: \file{controller/} and \file{bridge/} hold only the
files handed out, and the modules are the group's to write in its own repository. A testbench
whose modules are missing is skipped, never failed.

\textbf{Quartus Prime Lite and a DE0-CV board are recommended but not necessary.} No exercise and
no assessed element requires them; appendix~\ref{app:quartus} describes the flow for those who
want it, and for the group project's last session, where every group brings up its own node on
hardware handed out to them.

\section*{Where to read next}
The node built in part~\ref{part:can} is one half of a system. The other half, the C++ code that
talks to it over SPI, is written in another course, against the same register map and the same
protocol specification \textemdash{} and against the same protocol book, which is shared between
the two. Appendix~\ref{app:reading} describes it, the standards the design relates to, and four
directions to take the subject further.
